\documentclass[12pt,a4paper]{article}

\usepackage[top=1in,bottom=1in,left=1.1in,right=1.1in]{geometry}
\usepackage{amsmath,amssymb}
\usepackage{booktabs}
\usepackage{array}
\usepackage{enumitem}
\usepackage{parskip}
\usepackage{xcolor}

\definecolor{deepblue}{RGB}{0,51,102}

% ---- section style (no titlesec needed) ----
\makeatletter
\renewcommand{\section}{\@startsection{section}{1}{\z@}
  {-2.5ex plus -1ex minus -.2ex}{1.2ex plus .2ex}
  {\large\bfseries\color{deepblue}}}
\makeatother

\pagestyle{plain}

% ============================================================
\begin{document}
% ============================================================

% ---- TITLE ----
\begin{center}
  {\small Department of Computer Engineering \quad|\quad
   Pune Institute of Computer Technology (PICT)}\\[0.4em]
  \noindent\rule{\linewidth}{1pt}\\[0.5em]
  {\Large\bfseries\color{deepblue}Mini Project 2 --- Initial Project Report}\\[0.3em]
  {\normalsize\bfseries LAO Laboratory \quad|\quad Academic Year 2025--26}\\[0.3em]
  \noindent\rule{\linewidth}{1pt}\\[0.8em]
\end{center}

% ---- TEAM TABLE ----
\begin{center}
\begin{tabular}{lll}
\toprule
\textbf{Name} & \textbf{Roll Number} & \textbf{Responsibility} \\
\midrule
Naman Lodha     & 16015024039 & SRFT, Interpolative Decomp., CUR \\
Aditya Lokhande & 16015024040 & Experiments, Visualizations, Error Analysis \\
Aryan Lomte     & 16015024041 & Core Randomized SVD (Stage A \& B) \\
\bottomrule
\end{tabular}
\end{center}

\begin{center}
\textbf{Division:} A \quad \textbf{Batch:} A2 \quad
\textbf{Submission Date:} March 19, 2026
\end{center}

\noindent\rule{\linewidth}{0.6pt}

% ============================================================
\section{Project Title}
% ============================================================

\begin{center}
\large\textbf{Finding Structure with Randomness: Probabilistic Algorithms for
Constructing Approximate Matrix Decompositions}
\end{center}

\noindent\textit{Based on}: N. Halko, P.\,G. Martinsson, J.\,A. Tropp,
arXiv:0909.4061v2 [math.NA], 2010. (SIAM Review 53(2):217--288, 2011.)

% ============================================================
\section{Project Description}
% ============================================================

Classical matrix decompositions --- particularly the Singular Value Decomposition (SVD)
and QR factorization --- are fundamental to data analysis and scientific computing. For
large-scale matrices, the $\mathcal{O}(mn\min(m,n))$ cost of classical SVD becomes
prohibitive. This project implements and experimentally validates the
\textbf{randomized two-stage framework} for low-rank matrix approximation proposed by
Halko, Martinsson, and Tropp (2009).

The core idea is to replace expensive deterministic operations with random sampling:
a random test matrix $\Omega \in \mathbb{R}^{n \times l}$ computes a sketch $Y = A\Omega$
whose column space approximates the dominant subspace of $A$. After orthonormalization,
this basis $Q$ is used in a cheap deterministic Stage B to produce approximate SVDs at
a fraction of the classical cost. The method is underpinned by rigorous probabilistic
error bounds and beats classical algorithms in speed, robustness, and suitability for
parallel architectures.

% ============================================================
\section{Concept Used}
% ============================================================

\begin{itemize}[leftmargin=1.8em,itemsep=2pt]
  \item \textbf{Randomized Singular Value Decomposition (Randomized SVD)}
  \item \textbf{QR Factorization} --- stabilized Gram--Schmidt orthonormalization
  \item \textbf{Power / Subspace Iteration} --- improves accuracy for slowly-decaying spectra (Alg.~4.4)
  \item \textbf{Subsampled Randomized Fourier Transform (SRFT)} --- $\mathcal{O}(mn\log k)$ range finder (Alg.~4.5)
  \item \textbf{Interpolative Decomposition (ID)} --- column-subset low-rank representation (Sec.~5.2)
  \item \textbf{CUR Decomposition} --- row and column submatrix factorization (Sec.~2.1.4)
  \item \textbf{Johnson--Lindenstrauss Lemma \& Measure Concentration} --- theoretical error analysis
  \item \textbf{Principal Component Analysis (PCA)} --- application of randomized SVD
\end{itemize}

The central identity exploited is:
\[
A \;\approx\; Q\,Q^\top A, \qquad Q \in \mathbb{R}^{m \times l},\quad
Q^\top Q = I_l,\quad l = k + p,
\]
where $k$ is the target rank and $p \geq 0$ is the oversampling parameter.

% ============================================================
\section{Application Area}
% ============================================================

\begin{enumerate}[leftmargin=1.8em,itemsep=3pt]
  \item \textbf{Image Compression} --- truncated randomized SVD for image reconstruction
    at varying compression ratios
  \item \textbf{Dimensionality Reduction / PCA} --- fast approximate PCA on large datasets
  \item \textbf{Recommendation Systems} --- low-rank approximation of user--item matrices
  \item \textbf{Large-Scale Numerical Linear Algebra} --- pass-efficient algorithms
    for out-of-core matrices
\end{enumerate}

% ============================================================
\section{Expected Outcome}
% ============================================================

\textbf{What will be implemented in code:}
\begin{enumerate}[leftmargin=1.8em,itemsep=2pt]
  \item \textbf{Algorithm 4.1}: Basic Randomized Range Finder (Gaussian test matrix)
  \item \textbf{Algorithm 4.2}: Adaptive Range Finder (unknown rank, fixed-precision)
  \item \textbf{Algorithm 4.4}: Stabilized Power Iteration Range Finder ($q$ steps)
  \item \textbf{Algorithm 4.5}: SRFT-Based Fast Range Finder ($\mathcal{O}(mn\log l)$ cost)
  \item \textbf{Algorithm 5.1}: Direct SVD from basis $Q$ (Stage B)
  \item \textbf{Section 5.2}: Randomized Interpolative Decomposition
  \item \textbf{Section 2.1.4}: CUR Decomposition
\end{enumerate}

\textbf{Expected experimental results:}
\begin{enumerate}[leftmargin=1.8em,itemsep=3pt]
  \item Randomized SVD matches classical SVD accuracy:
    $\|A - \tilde{A}\|_F / \|A\|_F \approx 10^{-3}$ on well-conditioned matrices.
  \item $3$--$11\times$ wall-time speedup over \texttt{scipy.linalg.svd} at $m \geq 800$.
  \item \textbf{Theorem 1.1}: Empirical error converges to $\sigma_{k+1}$ as oversampling
    $p$ increases from 0 to 50.
  \item \textbf{Theorem 1.2}: Power iteration $q=2$ reduces error by $>10\times$ over $q=0$
    on flat-spectrum matrices.
  \item SRFT achieves near-identical accuracy to Gaussian with lower asymptotic cost.
  \item Image compression curves: compression ratio vs.\ reconstruction error at
    $k \in \{5, 20, 50, 100\}$.
\end{enumerate}

% ============================================================
\section{Dataset Used}
% ============================================================

\begin{center}
\begin{tabular}{lll}
\toprule
\textbf{Dataset} & \textbf{Source} & \textbf{Purpose} \\
\midrule
Synthetic low-rank matrices & Generated in-code (NumPy) & Benchmarking \& theorem verification \\
Synthetic image ($256{\times}256$) & Generated in-code (NumPy) & Image compression demo \\
MNIST (optional extension) & Kaggle (URL: kaggle.com/datasets/hojjatk/mnist-dataset) & Real-data PCA validation \\
\bottomrule
\end{tabular}
\end{center}

\noindent Synthetic matrices are constructed as
$A = U\,\mathrm{diag}(\sigma_1,\ldots,\sigma_r)\,V^\top + \eta E$,
where $\sigma_j = j^{-\alpha}$ controls spectral decay. No external data download
is required for the core experiments.

% ============================================================
\section{Tools and Libraries}
% ============================================================

\begin{center}
\begin{tabular}{>{\bfseries}ll}
\toprule
Library & Purpose \\
\midrule
NumPy & Matrix operations, random generation, linear algebra \\
SciPy (\texttt{linalg}, \texttt{fft}) & QR decomposition, SVD, FFT for SRFT \\
Matplotlib & Visualization, benchmark dashboards, experiment plots \\
\texttt{time} / \texttt{timeit} & Accurate wall-clock benchmarking \\
Python 3.12 & Implementation language \\
\bottomrule
\end{tabular}
\end{center}

% ============================================================
\section{Tentative Algorithm / Methodology}
% ============================================================

\textbf{Stage A: Randomized Range Finder (Algorithm 4.4)}
\begin{enumerate}[leftmargin=1.8em,itemsep=3pt]
  \item Draw $\Omega \in \mathbb{R}^{n \times l}$ with i.i.d.\ $\mathcal{N}(0,1)$ entries,
    $l = k + p$.
  \item Compute $Y = A\Omega \in \mathbb{R}^{m \times l}$.
  \item Apply $q$ steps of power iteration with QR stabilization:
    \[
    Y \;\leftarrow\; A\;\mathrm{qr}(A^\top\;\mathrm{qr}(Y)).
    \]
  \item Orthonormalize: $(Q, \_) = \mathrm{QR}(Y)$.
    Return $Q \in \mathbb{R}^{m \times l}$ with $A \approx QQ^\top A$.
\end{enumerate}

\textbf{Stage B: Direct SVD (Algorithm 5.1)}
\begin{enumerate}[leftmargin=1.8em,itemsep=3pt]
  \item Form $B = Q^\top A \in \mathbb{R}^{l \times n}$ (small matrix, $l \ll m$).
  \item Compute SVD: $B = \hat{U}\,\Sigma\,V^\top$.
  \item Set $U = Q\hat{U}$.
  \item Return truncated SVD: $U_k, \Sigma_k, V_k^\top$
\end{enumerate}

\textbf{SRFT Acceleration (Algorithm 4.5)}: Replace Gaussian $\Omega$ with
$\Omega = \sqrt{n/l}\,R\,F\,D$ (random sign flip $D$, FFT $F$, column subsample $R$).
Cost: $\mathcal{O}(mn\log l)$ vs.\ $\mathcal{O}(mnl)$.

\textbf{Experimental Validation Plan (5 Experiments):}
\begin{enumerate}[leftmargin=1.8em,itemsep=2pt]
  \item \textbf{Image Compression}: Ranks $k \in \{5,20,50,100\}$; measure
    $\|A-\tilde{A}\|_F/\|A\|_F$.
  \item \textbf{Oversampling Effect}: Plot error vs.\ $p \in [0,50]$; verify Theorem~1.1.
  \item \textbf{Power Iteration}: Error vs.\ $q \in \{0,1,2,3\}$ for three decay rates;
    verify Theorem~1.2.
  \item \textbf{PCA}: Explained variance comparison, classical vs.\ randomized.
  \item \textbf{Benchmark Dashboard}: Time, error, speedup vs.\ $m \in \{200,\ldots,1200\}$.
\end{enumerate}

% ============================================================
\section{Team Members with Name \& Roll Numbers}
% ============================================================

\begin{center}
\begin{tabular}{llll}
\toprule
\textbf{Name} & \textbf{Roll No.} & \textbf{Module File} & \textbf{Responsibility} \\
\midrule
Naman Lodha     & 16015024039 & \texttt{srft\_and\_decompositions.py}     & Alg.~4.5, ID (Sec.~5.2), CUR (Sec.~2.1.4) \\
Aditya Lokhande & 16015024040 & \texttt{experiments\_and\_applications.py} & All 5 experiments, plots \\
Aryan Lomte     & 16015024041 & \texttt{randomized\_svd\_core.py}          & Alg.~4.1, 4.2, 4.4, 5.1, benchmarking \\
\bottomrule
\end{tabular}
\end{center}

% ============================================================
\section{Division / Batch}
% ============================================================

\begin{center}
\begin{tabular}{ll}
\toprule
\textbf{Field} & \textbf{Detail} \\
\midrule
Department      & Computer Science and Business Systems \\
Division        & A \\
Batch           & A2 \\
Subject         & Linear Algebra and Optimization (LAO) \\
Academic Year   & 2025--26 \\
Submission Date & March 19, 2026 \\
\bottomrule
\end{tabular}
\end{center}

\vspace{1.5em}
\noindent\rule{\linewidth}{0.6pt}
\vspace{0.3em}

\noindent\small\textit{%
\textbf{Reference}: N. Halko, P.\,G. Martinsson, J.\,A. Tropp,
``Finding Structure with Randomness: Probabilistic Algorithms for Constructing
Approximate Matrix Decompositions,''
SIAM Review 53(2):217--288, 2011.
\texttt{arXiv:0909.4061v2}.
}

\end{document}
